\documentclass[aspectratio=169]{beamer}
\input{header}
\coursecode{JKSSB}

\title{CPU, ALU, Control Unit \& Registers}
\subtitle{Lesson 05 --- Functional Units and Registers}
\author{JKSSB Computer Awareness}
\date{\today}

\begin{document}
\frame{\titlepage}

\begin{frame}{Video Lesson 05}
  \centering
  \Large\href{https://youtu.be/b6w8uT4GBYQ}{Watch: CPU Explained: ALU, Control Unit \& Registers}
\end{frame}

\begin{frame}{The CPU and Its Functional Units}
  \begin{itemize}
    \item The \key{Central Processing Unit (CPU)} is often called the
      ``brain'' of the computer.
    \item It fetches, decodes, and executes instructions and works with main
      memory.
    \item Its main functional units are the Arithmetic Logic Unit (ALU),
      the Control Unit (CU), and the registers.
    \item Data and control signals move between these units over the
      internal buses.
  \end{itemize}
  \begin{block}{Memory hook}
    ALU computes, the Control Unit coordinates, and the registers hold the
    working data.
  \end{block}
\end{frame}

\begin{frame}{The Three Functional Units at a Glance}
  \begin{columns}[T]
    \column{0.32\textwidth}
      \begin{block}{ALU}
        Performs arithmetic and logical operations on binary data.
      \end{block}
    \column{0.32\textwidth}
      \begin{block}{Control Unit}
        Directs and sequences operations; issues control signals.
      \end{block}
    \column{0.32\textwidth}
      \begin{block}{Registers}
        Smallest, fastest storage, located inside the CPU.
      \end{block}
  \end{columns}
  \vspace{0.4cm}
  \muted{Keep the roles separate; most wrong answers swap the ALU and the
  Control Unit.}
\end{frame}

\begin{frame}{The Arithmetic Logic Unit (ALU)}
  \begin{itemize}
    \item Performs \key{arithmetic} operations: addition, subtraction,
      increment, decrement, multiplication, and division.
    \item Performs \key{logical} operations: AND, OR, NOT, XOR, shifts, and
      comparisons.
    \item Operates on \key{binary data} using logic-gate operations.
    \item Returns its result to a register or the accumulator; it does not
      permanently store data.
  \end{itemize}
  \begin{alertblock}{Exam cue}
    The ALU is the part that computes; its operations are carried out with
    logic gates on binary values.
  \end{alertblock}
\end{frame}

\begin{frame}{ALU Works on Binary Data with Logic Gates}
  \begin{itemize}
    \item Inside the CPU, numbers and characters are represented in binary.
    \item Logic gates such as AND, OR, NOT, and XOR implement the ALU's
      operations.
    \item Example: adding 1 and 1 works on the binary values, and the result
      is placed where it can be reused.
    \item \key{ALU} stands for Arithmetic Logic Unit.
  \end{itemize}
  \begin{block}{Remember}
    Arithmetic and logic are the two halves of the ALU's name; both run on
    binary data.
  \end{block}
\end{frame}

\begin{frame}{The Control Unit (CU)}
  \begin{itemize}
    \item \key{Directs and sequences} the operations of the computer.
    \item Controls the \key{flow of data} between the ALU, registers, and
      memory.
    \item Decodes instructions and \key{issues control signals} that time and
      trigger the other units.
    \item It does \bad{not} perform arithmetic or logical calculations.
  \end{itemize}
  \begin{block}{Role in one line}
    The Control Unit is the conductor: it tells every part what to do and
    when, but it does not itself compute.
  \end{block}
\end{frame}

\begin{frame}{Control Unit vs ALU: Do Not Swap the Roles}
  \begin{center}
    \begin{tabular}{ll}
      \toprule
      \textbf{Unit} & \textbf{Primary job} \\
      \midrule
      Arithmetic Logic Unit (ALU) & Performs arithmetic and logic operations \\
      Control Unit (CU) & Directs operations and issues control signals \\
      \bottomrule
    \end{tabular}
  \end{center}
  \vspace{0.35cm}
  \begin{alertblock}{Trap alert}
    The Control Unit controls and coordinates; it does \bad{not} add or
    subtract. Saying ``CU computes'' or ``ALU controls'' is wrong.
  \end{alertblock}
\end{frame}

\begin{frame}{Registers: Fastest Storage Inside the CPU}
  \begin{itemize}
    \item \key{Registers} are the \key{smallest and fastest} storage in a
      computer.
    \item They sit inside the CPU and are faster than cache and RAM.
    \item They hold instructions, addresses, data, and intermediate results
      temporarily while work is in progress.
    \item Their capacity is very small, so they are used only for the
      values needed right now.
  \end{itemize}
  \begin{block}{Speed order}
    Registers are fastest, then cache, then main memory (RAM), then
    secondary storage.
  \end{block}
\end{frame}

\begin{frame}{Key Registers and Their Purposes}
  \begin{center}
    \begin{tabular}{ll}
      \toprule
      \textbf{Register} & \textbf{Holds} \\
      \midrule
      Program Counter (PC / IP) & Address of the next instruction to fetch \\
      Instruction Register (IR) & Instruction currently being decoded/executed \\
      Memory Address Register (MAR) & Address of the memory location to access \\
      Memory Data Register (MDR) & Data being transferred to or from memory \\
      Accumulator (ACC) & Intermediate arithmetic/logic results \\
      General-purpose registers & Temporary operands and results \\
      Status / flags register & Condition flags (zero, carry, sign, overflow) \\
      \bottomrule
    \end{tabular}
  \end{center}
  \vspace{0.2cm}
  \muted{PC/IP = Program Counter or Instruction Pointer.}
\end{frame}

\begin{frame}{The Program Counter: A Common Trap}
  \begin{itemize}
    \item The \key{Program Counter (PC)} holds the address of the
      \good{next} instruction to be fetched.
    \item It does \bad{not} hold the address of the instruction currently
      being executed.
    \item The \key{Instruction Register (IR)} holds the instruction that is
      currently being decoded or executed.
  \end{itemize}
  \begin{alertblock}{Trap alert}
    Near-neighbour confusion: PC = address of the \good{next}
    instruction; IR = the \good{current} instruction itself.
  \end{alertblock}
\end{frame}

\begin{frame}{Clock Speed, Cores, and Threads}
  \begin{itemize}
    \item \key{Clock speed} is measured in hertz (Hz); 1 GHz is one billion
      cycles per second.
    \item A \key{clock cycle} is one tick of the system clock; different
      instructions may take different numbers of cycles.
    \item A \key{core} is an independent processing unit inside the CPU; a
      quad-core chip has four cores.
    \item A \key{thread} is a stream of instructions; multiple threads can
      share or run on the available cores.
  \end{itemize}
  \begin{block}{Office desktop example}
    A quad-core 2.4 GHz machine has four cores, each running at about 2.4
    billion clock cycles per second.
  \end{block}
\end{frame}

\begin{frame}{Lesson 05: Recall}
  \begin{itemize}
    \item The CPU has the ALU, the Control Unit, and the registers.
    \item The ALU performs arithmetic and logic on binary data using logic
      gates; the Control Unit only directs and issues control signals.
    \item Registers are the fastest, smallest storage, inside the CPU.
    \item PC holds the address of the \good{next} instruction; IR holds the
      \good{current} instruction; MAR holds an address; MDR holds data; the
      accumulator holds intermediate results.
    \item Clock speed is in Hz/GHz; cores and threads describe parallel
      execution.
  \end{itemize}
\end{frame}
\end{document}
